\documentclass[11pt]{article}
\usepackage[margin=1in,landscape]{geometry}
\usepackage[T1]{fontenc}
\usepackage[utf8]{inputenc}
\usepackage[brazilian]{babel}

\usepackage{array}
\usepackage{booktabs}
\usepackage{xcolor}
\usepackage{listings}
\usepackage{hyperref}
\usepackage{comment}


\setlength{\parindent}{4pt}
\setlength{\parskip}{1\baselineskip}

\lstdefinestyle{csvstyle}{
  basicstyle=\ttfamily\small,
  backgroundcolor=\color{gray!10},
  frame=single,
  breaklines=true,
  columns=fullflexible
}

\lstdefinestyle{cppstyle}{
  language=MATLAB,
  basicstyle=\ttfamily\small,
  keywordstyle=\color{blue!70!black}\bfseries,
  commentstyle=\color{green!40!black},
  stringstyle=\color{orange!60!black},
  backgroundcolor=\color{gray!8},
  frame=single,
  numbers=left,
  numberstyle=\tiny\color{gray},
  breaklines=true,
  columns=fullflexible,
  tabsize=2,
  showstringspaces=false
}

\title{Resultados de TP: Simulação de um sistema termodinâmico turbina-compressor-queimador}
\author{Lucas Ferreira Simões (n°2026020102)\\Universidade Federal de Minas Gerais}
\date{\today}

\begin{document}
\maketitle

\section{Introdução}

O trabalho é sobre um sistema compressor-queimador-turbina, onde o compressor começa ativado (gastando Ec = 1000kW), pegando ar do ambiente (t1 = 25°C, p = 101kPa) e enviando (vazão w = 10kg/s) para dentro do sistema (t2 = 250°C, p2 = 300kPa) a caminho do queimador, que aumenta a temperatura desse ar (t (ou t3) = 800°C, p3 = p2 = 300kPa) e o injeta na turbina que gera movimento para realimentar o compressor e sobrar energia para qualquer outra carga estabilizada (Es(carga) = 1500kW, Et(turbina) = 2400kW). Pelos valores iniciais, a ideia alterar os valores das variáveis para convergir a um ponto de equilíbrio entre todas elas, isto é, onde a turbina ofereça a energia exata para o compressor enviar a quantidade exata de ar ao queimador para manter o circuito estável e funcionando. Para isso, resolvemos um sistema de equações cujas variáveis são as componentes dessas alterações no circuito.

Nesse trabalho uso o método de Newton-Raphson para várias variáveis, resolvendo sistemas usando fatoração LU e matriz jacobiana para aproximar o sistema de funções simultâneas. Como o professor não pediu um relatório completo, não vou me estender aqui. O pedido foi que eu apresentasse os resultados dados por esses valores iniciais.

O trabalho foi feito e testado em Octave, depois compilado e testado em MATLAB e apresentou resultados semelhantes nos dois, funcionando perfeitamente. Os únicos warnings que o compilador apresentou foram sobre alocar grande espaço para os vetores (ao invés de dinamicamente) para melhorar a velocidade do código, mas não é necessário e nem perceptível nesse trabalho.


\section{Resultados}

Os resultados do método de Newton-Raphson são apresentados por iteração. O erro e o vetor de correção $\Delta$ ocupam colunas separadas.

\begin{table}[htbp]
\centering
\caption{Resultados por iteração do método de Newton-Raphson}
\scriptsize
\setlength{\tabcolsep}{2.4pt}
\renewcommand{\arraystretch}{3}
\begin{tabular}{@{}r r r r r r r r r m{5.5cm}@{}}
\toprule
$k$ & $w$ (kg/s) & $p$ (kPa) & $E_c$ (kW) & $t_2$ ($^\circ$C) & $E_s$ (kW) & $t_3$ ($^\circ$C) & $E_t$ (kW) & Erro & $\Delta$ \\
\midrule
0 & 10.0000000000 & 300.0000000000 & 1000.0000000000 & 250.0000000000 & 1500.0000000000 & 800.0000000000 & 2400.0000000000 & 1.00e+00 &  \\
1 & 10.9093592993 & 352.1724245256 & 1507.4046840965 & 150.8893850962 & 1570.6849821317 & 877.5736527627 & 3078.0896662283 & 6.78e+02 & \shortstack[l]{9.09e-01, 5.22e+01, 5.07e+02, -9.91e+01 \\ 7.07e+01, 7.76e+01, 6.78e+02} \\
2 & 10.7677146962 & 354.8356153068 & 1529.9712246322 & 162.7936565012 & 1597.7452527609 & 884.1853514664 & 3127.7164773931 & 4.96e+01 & \shortstack[l]{-1.42e-01, 2.66e+00, 2.26e+01, 1.19e+01 \\ 2.71e+01, 6.61e+00, 4.96e+01} \\
3 & 10.7652956155 & 354.8543972012 & 1530.0687937452 & 162.9900061562 & 1598.5137013795 & 884.4741478838 & 3128.5824951247 & 8.66e-01 & \shortstack[l]{-2.42e-03, 1.88e-02, 9.76e-02, 1.96e-01 \\ 7.68e-01, 2.89e-01, 8.66e-01} \\
4 & 10.7652949442 & 354.8544012565 & 1530.0688087663 & 162.9900602383 & 1598.5140917183 & 884.4742677338 & 3128.5829004846 & 4.05e-04 & \shortstack[l]{-6.71e-07, 4.06e-06, 1.50e-05, 5.41e-05 \\ 3.90e-04, 1.20e-04, 4.05e-04} \\
5 & 10.7652949442 & 354.8544012565 & 1530.0688087663 & 162.9900602383 & 1598.5140917183 & 884.4742677339 & 3128.5829004846 & 3.75e-11 & \shortstack[l]{-2.82e-14, -6.00e-13, -1.84e-12, 3.58e-12 \\ 3.75e-11, 9.65e-12, 3.59e-11} \\
\bottomrule
\end{tabular}
\end{table}

\begin{center}
A ordem dos componentes de $\Delta$ é $(w,p,E_c,t_2,E_s,t_3,E_t)$. A iteração $k=0$ não possui vetor de correção.
\end{center}

Pelos resultados, as variáveis convergiram para aproximadamente $w = 10{,}7653\,\mathrm{kg/s}$, $p = 354{,}8544\,\mathrm{kPa}$, $E_c = 1530{,}0688\,\mathrm{kW}$, $t_2 = 162{,}9901\,^{\circ}\mathrm{C}$, $E_s = 1598{,}5141\,\mathrm{kW}$, $t_3 = 884{,}4743\,^{\circ}\mathrm{C}$ e $E_t = 3128{,}5829\,\mathrm{kW}$ (os valores com 10 casas decimais estão na tabela). Portanto, o sistema estaria em equilíbrio aproximadamente nesses valores. O erro máximo obtido foi de $3{,}75 \times 10^{-11}$, abaixo do limite de $1 \times 10^{-10}$ indicado no enunciado. Os valores dos polinômios de ajuste $\Delta$ também constam da tabela.

Notei que o valor de t3 foi o último a convergir. Isso acontece porque ele está presente em várias equações, e elas precisam todas começarem a estabilizar para se estabilizar também. Testei o valor inicial dele em 884 e ele voltou para 870, para daí então convergir normalmente. Em inicial=870, ele foi pra 871. Geometricamente, parece que a subida dele foi mais útil quando começou de longe do que quando começou de perto.
 
\end{document}